% ===================================================================== PREAMBULE COMMUN — Carnet d'entraînement Fiche 9
\documentclass[11pt,a4paper]{article}
\usepackage[utf8]{inputenc}
\usepackage[T1]{fontenc}
\usepackage{lmodern}
\usepackage[french]{babel}
\usepackage{amsmath,amssymb,mathtools}
\usepackage{icomma}
\usepackage{siunitx}
\sisetup{output-decimal-marker={,},per-mode=symbol}
\DeclareSIUnit{\molar}{M}
\DeclareSIUnit{\Molar}{mol\,L^{-1}}
\usepackage[version=4]{mhchem}
\usepackage{xcolor}
\usepackage{colortbl}
\usepackage{tikz}
\usepackage{pgfplots}
\usepackage[most]{tcolorbox}
\usepackage{enumitem}
\usepackage{array}
\usepackage{booktabs}
\usepackage{multicol}
\usepackage{fancyhdr}
\usepackage[hidelinks]{hyperref}
\usetikzlibrary{arrows.meta,positioning,calc,decorations.pathmorphing,
  decorations.markings,angles,quotes,patterns,backgrounds,
  shapes.geometric,shapes.misc,fadings,intersections,fit}
\pgfplotsset{compat=1.18}
\usepackage[a4paper,margin=2.2cm,headheight=26pt,headsep=10pt]{geometry}

% ---------- Palette ----------
\definecolor{primary}{RGB}{150,98,162}
\definecolor{primarydark}{RGB}{92,52,110}
\definecolor{defcol}{RGB}{0,114,178}
\definecolor{thmcol}{RGB}{0,140,120}
\definecolor{formulacol}{RGB}{198,93,9}
\definecolor{remarkcol}{RGB}{120,94,200}
\definecolor{attncol}{RGB}{200,40,55}
\definecolor{excol}{RGB}{0,135,90}
\definecolor{methcol}{RGB}{90,90,100}
\definecolor{cationcol}{RGB}{0,90,190}
\definecolor{anioncol}{RGB}{200,60,55}
\definecolor{cristalcol}{RGB}{110,105,120}
\definecolor{eaucol}{RGB}{120,180,220}
\definecolor{solcol}{RGB}{0,135,90}
\definecolor{kpscol}{RGB}{198,93,9}
\definecolor{qcol}{RGB}{120,94,200}
\definecolor{precipcol}{RGB}{110,70,40}
\definecolor{exercol}{RGB}{150,98,162}
\definecolor{corrcol}{RGB}{0,120,100}

% ---------- Macros ----------
\newcommand{\Kps}{K_{\mathrm{ps}}}
\newcommand{\Ce}[1]{\left[\,\ce{#1}\,\right]}

% ---------- Style des boîtes ----------
\tcbset{boxbase/.style={enhanced,breakable,boxrule=0.9pt,arc=2.5pt,
   left=3mm,right=3mm,top=2.3mm,bottom=2.3mm,
   fonttitle=\bfseries,coltitle=white,toptitle=0.6mm,bottomtitle=0.6mm}}

\newtcolorbox[auto counter]{definition}[1][]{%
  boxbase,colback=defcol!5,colframe=defcol,
  title={\textsc{Définition}~\thetcbcounter\ \ifx\relax#1\relax\relax\else\textmd{--- \itshape #1}\fi}}
\newtcolorbox[auto counter]{theoreme}[1][]{%
  boxbase,colback=thmcol!6,colframe=thmcol,
  title={\textsc{Loi / Propriété}~\thetcbcounter\ \ifx\relax#1\relax\relax\else\textmd{--- \itshape #1}\fi}}
\newtcolorbox[auto counter]{formule}[1][]{%
  boxbase,colback=formulacol!6,colframe=formulacol,
  title={\textsc{Formule}~\thetcbcounter\ \ifx\relax#1\relax\relax\else\textmd{--- \itshape #1}\fi}}
\newtcolorbox{remarque}[1][]{%
  boxbase,colback=remarkcol!6,colframe=remarkcol,
  title={\textsc{Remarque}\ifx\relax#1\relax\relax\else\textmd{ : \itshape #1}\fi}}
\newtcolorbox{attention}[1][]{%
  boxbase,colback=attncol!6,colframe=attncol,
  title={\textsc{Attention}\ifx\relax#1\relax\relax\else\textmd{ : \itshape #1}\fi}}
\newtcolorbox{exemple}[1][]{%
  boxbase,colback=excol!5,colframe=excol,
  title={\textsc{Exemple}\ifx\relax#1\relax\relax\else\textmd{ : \itshape #1}\fi}}
\newtcolorbox{methode}[1][]{%
  boxbase,colback=methcol!7,colframe=methcol,sharp corners,
  title={\textsc{Méthode}\ifx\relax#1\relax\relax\else\textmd{ : \itshape #1}\fi}}
\newtcolorbox{astuce}[1][]{%
  boxbase,colback=primary!7,colframe=primary,
  title={\textsc{Astuce}\ifx\relax#1\relax\relax\else\textmd{ : \itshape #1}\fi}}

\newtcolorbox[auto counter]{exercice}[1][]{%
  boxbase,colback=exercol!4,colframe=exercol,
  title={\textsc{Exercice}~\thetcbcounter\ \ifx\relax#1\relax\relax\else\textmd{--- \itshape #1}\fi}}
\newtcolorbox[auto counter]{correction}[1][]{%
  boxbase,colback=corrcol!4,colframe=corrcol,
  title={\textsc{Corrigé}~\thetcbcounter\ \ifx\relax#1\relax\relax\else\textmd{--- \itshape #1}\fi}}

% ---------- Environnement « où : » ----------
\newenvironment{ou}{%
  \par\smallskip\noindent\textbf{où :}\nopagebreak
  \begin{itemize}[leftmargin=1.8em,topsep=2pt,itemsep=1.5pt,parsep=0pt]}%
  {\end{itemize}}

% ---------- Styles TikZ ----------
\tikzset{
  cat/.style={circle,fill=cationcol,draw=cationcol!50!black,inner sep=0pt,
              minimum size=5mm,font=\bfseries\scriptsize,text=white},
  an/.style={circle,fill=anioncol,draw=anioncol!50!black,inner sep=0pt,
             minimum size=5mm,font=\bfseries\scriptsize,text=white},
  crx/.style={circle,fill=cristalcol,draw=black!50,inner sep=0pt,
              minimum size=5mm,font=\bfseries\scriptsize,text=white},
  ptn/.style={circle,fill=black,inner sep=1.1pt}}

% ---------- En-têtes ----------
\pagestyle{fancy}
\fancyhf{}
\fancyhead[L]{\small\textcolor{primarydark}{\textbf{Fiche 9} — Solubilité et produit de solubilité}}
\fancyhead[R]{\small\textcolor{primarydark}{\textsc{Carnet d'entraînement}}}
\fancyfoot[C]{\small\thepage}
\renewcommand{\headrulewidth}{0.8pt}
\renewcommand{\headrule}{\hbox to\headwidth{\color{primary}\leaders\hrule height \headrulewidth\hfill}}
\usepackage{titlesec}
\titleformat{\section}{\Large\bfseries\color{primarydark}}{\thesection}{0.6em}{}
\titleformat{\subsection}{\large\bfseries\color{primary}}{\thesubsection}{0.6em}{}

% ---------- Bandeau de titre ----------
% #1 = thème/étiquette   #2 = titre principal   #3 = sous-titre
\newcommand{\bandeau}[3]{%
\begin{center}
\begin{tikzpicture}
\node[rectangle,rounded corners=7pt,fill=primary,text=white,
      minimum width=15.6cm,minimum height=1.95cm,align=center,inner sep=8pt]
  {{\large\bfseries #1}\\[3pt]{\Large\bfseries #2}\\[3pt]{\small #3}};
\end{tikzpicture}
\end{center}
\vspace{2mm}}
\begin{document}
\bandeau{Thème 5 — Influence du pH}%
{Corrigé — Niveau Difficile}%
{Réponses détaillées et justifications}

\begin{correction}[séparer le fer du magnésium]
\begin{enumerate}[label=\textbf{\alph*)},leftmargin=2.2em,topsep=2pt,itemsep=2pt]
  \item \ce{Fe(OH)3} : $[\ce{OH-}]=\sqrt[3]{\num{2{,}0e-39}/\num{1{,}0e-2}}=\sqrt[3]{\num{2{,}0e-37}}
        \approx\qty{5{,}9e-13}{\molar}$, $\mathrm{pH}\approx\boxed{1{,}8}$.\\
        \ce{Mg(OH)2} : $[\ce{OH-}]=\sqrt{\num{9{,}0e-12}/\num{1{,}0e-2}}=\sqrt{\num{9{,}0e-10}}
        =\qty{3{,}0e-5}{\molar}$, $\mathrm{pH}\approx\boxed{9{,}5}$.
  \item Entre ces deux seuils : pour $1{,}8\lesssim\mathrm{pH}\lesssim 9{,}5$, le fer est précipité mais
        le magnésium reste dissous. Un choix confortable est \textbf{$\mathrm{pH}\approx 4$ à $8$}
        (sur la figure, la courbe \ce{Fe(OH)3} est déjà très basse tandis que \ce{Mg(OH)2} est encore
        au plafond).
  \item À $\mathrm{pH}=7$ : $[\ce{OH-}]=10^{7-14}=\qty{1{,}0e-7}{\molar}$, donc
        $[\ce{Fe^3+}]=\dfrac{\Kps}{[\ce{OH-}]^{3}}=\dfrac{\num{2{,}0e-39}}{(\num{1{,}0e-7})^{3}}
        =\dfrac{\num{2{,}0e-39}}{\num{1{,}0e-21}}=\boxed{\qty{2{,}0e-18}{\molar}}$ : le fer est
        \emph{totalement} éliminé.
  \item Pour le magnésium à $\mathrm{pH}=7$ : $Q=[\ce{Mg^2+}][\ce{OH-}]^{2}
        =(\num{1{,}0e-2})(\num{1{,}0e-7})^{2}=\num{1{,}0e-16}$. Comme $Q=\num{1{,}0e-16}<\Kps=\num{9{,}0e-12}$,
        \ce{Mg(OH)2} ne précipite pas : le magnésium \textbf{reste en solution}. La séparation est donc
        \textbf{réalisable} à pH neutre.
\end{enumerate}
\end{correction}

\begin{correction}[prédire l'effet d'une acidification]
\begin{enumerate}[label=\textbf{\alph*)},leftmargin=2.2em,topsep=2pt,itemsep=2pt]
  \item \begin{itemize}[leftmargin=1.4em,topsep=1pt,itemsep=1pt]
    \item \ce{AgCl} : \textbf{ne change pas} — \emph{indifférent} (\ce{Cl-}, acide fort).
    \item \ce{CaCO3} : \textbf{augmente} — \emph{sel d'acide faible} (\ce{CO3^2-} consommé par \ce{H3O+}).
    \item \ce{Fe(OH)3} : \textbf{augmente} — \emph{hydroxyde} (baisser le pH, c'est retirer des \ce{OH-}).
    \item \ce{BaSO4} : \textbf{ne change pas} — \emph{indifférent} (\ce{SO4^2-}, traité comme acide fort).
    \item \ce{FeS} : \textbf{augmente} — \emph{sel d'acide faible} (\ce{S^2-} consommé par \ce{H3O+}).
  \end{itemize}
  \item $s=\dfrac{\Kps}{[\ce{OH-}]^{3}}=\Kps\cdot 10^{\,3(14-\mathrm{pH})}$. Donc
        \[ \dfrac{s(\mathrm{pH}=2)}{s(\mathrm{pH}=7)}=\dfrac{10^{3(14-2)}}{10^{3(14-7)}}
        =\dfrac{10^{36}}{10^{21}}=\boxed{10^{15}}. \]
        La solubilité de \ce{Fe(OH)3} est multipliée par $10^{15}$ : voilà pourquoi le fer(III), quasi
        insoluble à pH neutre, se redissout complètement en milieu très acide.
\end{enumerate}
\end{correction}

\begin{correction}[ion commun \emph{ou} pH : le même geste]
\begin{enumerate}[label=\textbf{\alph*)},leftmargin=2.2em,topsep=2pt,itemsep=2pt]
  \item $Q=[\ce{Mg^2+}][\ce{OH-}]^{2}=(\num{1{,}0e-3})(\num{1{,}0e-3})^{2}=\num{1{,}0e-9}$. Comme
        $Q=\num{1{,}0e-9}>\Kps=\num{9{,}0e-12}$, il y a \textbf{précipitation} de \ce{Mg(OH)2}.
  \item $\mathrm{pH}=14+\log[\ce{OH-}]=14+\log(\num{1{,}0e-3})=14-3=\boxed{11}$ (solution basique).
  \item C'est \textbf{les deux à la fois} : \ce{OH-} est l'ion commun de \ce{Mg(OH)2} (vision « ion
        commun », Thème 4) \emph{et} augmenter $[\ce{OH-}]$ signifie augmenter le pH (vision « pH »,
        Thème 5). Pour un \textbf{hydroxyde}, l'anion \emph{est} \ce{OH-} : « ajouter l'ion commun » et
        « monter le pH » sont un seul et même geste, et prédisent tous deux une baisse de solubilité
        (précipitation). Les deux thèmes se rejoignent ici.
\end{enumerate}
\end{correction}

\end{document}
